\begingroup
\renewcommand{\arraystretch}{1.45}
\fancyhead[L]{\small DETERMINATION OF ATTRACTION}
\Needspace{30\baselineskip}
\AttractionArticle{12}

Having shown how to determine completely the twelve quantities $G$, $G'$, $G''$, $\alpha$, $\alpha'$, $\alpha''$, $\beta$. $\beta'$, $\beta''$, $\gamma$, $\gamma'$, $\gamma''$, we proceed to calculate the differential $\mathrm dE$. Let us set
\[
\begin{aligned}
&t=\gamma+\gamma'\cos T+\gamma''\sin T\quad
\mathinner{\ldotp\;\ldotp\;\ldotp\;\ldotp\;\ldotp\;\ldotp\;\ldotp\;\ldotp}\quad[20]
\end{aligned}
\]
so that
\[
\begin{aligned}
&t\cos E=\alpha+\alpha'\cos T+\alpha''\sin T\quad
\mathinner{\ldotp\;\ldotp\;\ldotp\;\ldotp\;\ldotp\;\ldotp\;\ldotp}\quad[21]
\end{aligned}
\]
\[
\begin{aligned}
&t\sin E=\beta+\beta'\cos T+\beta''\sin T\quad
\mathinner{\ldotp\;\ldotp\;\ldotp\;\ldotp\;\ldotp\;\ldotp\;\ldotp}\quad[22]
\end{aligned}
\]
\Needspace{27\baselineskip}
Hence we deduce
\[
\begin{aligned}
&t\,\mathrm dE=\cos E\,\mathrm d.\,t\sin E-\sin E\,\mathrm d.\,t\cos E\\[8pt]
&=\cos E(\beta''\cos T-\beta'\sin T)\mathrm dT
-\sin E(\alpha''\cos T-\alpha'\sin T)\mathrm dT
\end{aligned}
\]
and therefore
\[
\begin{aligned}
&tt\,\mathrm dE=(\alpha\beta''-\alpha''\beta)\cos T\,\mathrm dT
+(\alpha'\beta-\beta'\alpha)\sin T\,\mathrm dT
\\[6pt]
&\qquad{}+(\alpha'\beta''-\beta'\alpha'')\mathrm dT\\[8pt]
&=\varepsilon\gamma'\cos T\,\mathrm dT
+\varepsilon\gamma''\sin T\,\mathrm dT
+\varepsilon\gamma\,\mathrm dT=\varepsilon t\,\mathrm dT
\end{aligned}
\]
or
\[
\begin{aligned}
&t\,\mathrm dE{}=\varepsilon\,\mathrm dT\quad
\mathinner{\ldotp\;\ldotp\;\ldotp\;\ldotp\;\ldotp\;\ldotp\;\ldotp\;\ldotp\;\ldotp\;\ldotp\;\ldotp\;\ldotp\;\ldotp\;\ldotp\;\ldotp\;\ldotp\;\ldotp\;\ldotp\;\ldotp\;\ldotp\;\ldotp}\quad[23]
\end{aligned}
\]
\Needspace{18\baselineskip}
\par It should be observed that the quantity $t$ is by its nature always positive if the coefficient $\gamma$ is positive, or always negative if $\gamma$ is negative. For, since
$(\gamma'\cos T+\gamma''\sin T)^2+(\gamma''\cos T-\gamma'\sin T)^2=\gamma'\gamma'+\gamma''\gamma''=\gamma\gamma-1$, the expression
\noindent\SourceLeaf{358}
$\gamma'\cos T+\gamma''\sin T$ will always, regardless of its sign, be less than $\gamma$. Hence we conclude that, whenever $\varepsilon\gamma$ is a positive quantity, the variables $E$ and $T$ always increase together; but whenever $\varepsilon\gamma$ is a negative quantity, one variable must always decrease while the other increases.

\Needspace{28\baselineskip}
\AttractionArticle{13}

The relation between the variables $E$ and $T$ is made still clearer by the following reasoning. Setting $\surd{(\gamma\gamma-1)}=\delta$, so that $\delta\delta=\alpha\alpha+\beta\beta=\gamma'\gamma'+\gamma''\gamma''$, we deduce from equations 20, 21, 22
\[
\begin{aligned}
&t(\delta+\alpha\cos E+\beta\sin E)\\[8pt]
&\quad=\gamma\delta+\alpha\alpha+\beta\beta
+(\gamma'\delta+\alpha\alpha'+\beta\beta')\cos T
+(\gamma''\delta+\alpha\alpha''+\beta\beta'')\sin T\\[8pt]
&\quad=(\gamma+\delta)(\delta+\gamma'\cos T+\gamma''\sin T)
\end{aligned}
\]
\Needspace{30\baselineskip}
\par Similarly, from equations 21, 22 it follows that
\[
\begin{aligned}
&t(\alpha\sin E-\beta\cos E)=\varepsilon(\gamma'\sin T-\gamma''\cos T)
\end{aligned}
\]
On setting
\[
\begin{aligned}
&\frac{\alpha}{\delta}=\cos L,\quad
\frac{\beta}{\delta}=\sin L,\quad
\frac{\gamma'}{\delta}=\cos M,\quad
\frac{\gamma''}{\delta}=\sin M
\end{aligned}
\]
these equations take the following form:
\[
\begin{aligned}
&t(1+\cos(E-L))=(\gamma+\delta)(1+\cos(T-M))
\end{aligned}
\]
\[
\begin{aligned}
&t\sin(E-L)=\varepsilon\sin(T-M)
\end{aligned}
\]
whence, by division, since $(\gamma+\delta)(\gamma-\delta)=1$,
\[
\begin{aligned}
&\mathrm{tang}\,\frac12(E-L)=\varepsilon(\gamma-\delta)\mathrm{tang}\,\frac12(T-M)
\end{aligned}
\]
\[
\begin{aligned}
&\mathrm{tang}\,\frac12(T-M)=\varepsilon(\gamma+\delta)\mathrm{tang}\,\frac12(E-L)
\end{aligned}
\]
\Needspace{18\baselineskip}
\par Hence not only is the same conclusion derived as the one to which we were led at the end of the preceding article, but it is also evident that, if the value of $E$ increases by 360 degrees, the value of $T$ either increases or decreases by the same amount, according as $\varepsilon\gamma$ is a positive or a negative quantity. Furthermore, on setting $\delta=\mathrm{tang}\,N$, $\gamma=\sec N$, it will evidently follow that
\[
\begin{aligned}
&\gamma-\delta=\mathrm{tang}(45^0-\frac12N),\quad
\gamma+\delta=\mathrm{tang}(45^0+\frac12N)
\end{aligned}
\]
\SourceLeaf{359}
\clearpage
\endgroup
